\section{Введение}

\begin{frame}
    \frametitle{Актуальность}
    \small
    \begin{itemize}
        \item Межпланетные перелёты с ЭРДУ формируются в результате длительного управляемого движения.
        \item На ранних этапах баллистического проектирования требуется многократно оценивать даты старта, продолжительность перелёта, расход рабочего тела и параметры управления.
        \item Непрямой метод требует начального вектора сопряжённых переменных, чувствительного к выбранной ветви решения.
    \end{itemize}

    \vspace{0.2cm}
    \[
    \begin{array}{c}
    t_0,\ \text{целевая орбита}
    \Rightarrow
    \text{выбор ветви}
    \Rightarrow
    T,\ J,\ \mathbf z_0\\
    \Rightarrow
    \text{начальный вектор}
    \Rightarrow
    \text{уточнение КЗ}.
    \end{array}
    \]
\end{frame}

\begin{frame}[t]
    \frametitle{Проблема проектирования класса перелётов}
    \begin{columns}[T,onlytextwidth]
        \begin{column}{0.56\textwidth}
            \centering
            \includegraphics[width=\linewidth,height=0.54\textheight,keepaspectratio]{Startdate}
        \end{column}
        \begin{column}{0.40\textwidth}
            \scriptsize
            \textbf{Одна краевая задача}
            \begin{itemize}
                \item фиксированы дата старта и целевая орбита;
                \item начальный вектор сопряжённых переменных подбирается индивидуально;
                \item решение уточняется численно.
            \end{itemize}

            \textbf{Класс проектных задач}
            \begin{itemize}
                \item дата старта меняет начальную синодическую фазу;
                \item возникают разные ветви решений;
                \item требуется согласованный способ выбора ветви и начального вектора.
            \end{itemize}
        \end{column}
    \end{columns}
    \vspace{-0.08cm}
    \tiny
    \[
    \begin{array}{c}
    t_0 \longrightarrow \Phi_0(t_0),\qquad \Phi_0=\mathrm{const}\\
    \text{задаёт условие выбора ветви.}
    \end{array}
    \]
\end{frame}

\begin{frame}
    \frametitle{Степень разработанности и место работы}
    \small
    \begin{columns}[T,onlytextwidth]
        \begin{column}{0.48\textwidth}
            \textbf{Конечные формулы, аналитические и усреднённые подходы}

            \medskip
            Лебедев, Белецкий, Гродзовский;\\
            da Silva Fernandes, Sessin, Carvalho, de Moraes;\\
            Баранов.

            \medskip
            \textbf{Интерпретируемые формулы по численным данным}

            \medskip
            Cranmer; символьная регрессия как инструмент получения конечных формул.
        \end{column}
        \begin{column}{0.48\textwidth}
            \textbf{Непрямые методы, продолжение, численные краевые задачи}

            \medskip
            Охоцимский, Энеев, Константинов, Петухов, Epenoy.

            \smallskip
            \textbf{Место работы}
            \[
            \begin{array}{c}
            \text{модельное многообразие}
            \Downarrow\\
            \text{конечные формулы}
            \Downarrow\\
            \text{начальный вектор}\\
            \text{для уточнения КЗ}.
            \end{array}
            \]
        \end{column}
    \end{columns}
\end{frame}

\begin{frame}
    \frametitle{Цель и задачи}
    \small
    \textbf{Цель:} разработать и апробировать численно-аналитический метод построения начальных приближений для энергетически оптимальных межпланетных перелётов с непрерывной тягой на основе аппроксимации модельного многообразия решений задачи встречи.

    \medskip
    Для достижения цели решаются задачи:
    \begin{enumerate}
        \item построить регулярное модельное многообразие решений задачи встречи;
        \item получить конечные формулы для характеристик экстремалей и начального вектора сопряжённых переменных;
        \item разработать линейные \(R\)-преобразования для перехода к эллиптической и пространственной конечной орбите;
        \item построить алгоритм выбора начального вектора и проверить его на прикладных классах перелётов.
    \end{enumerate}
\end{frame}

\begin{frame}
    \frametitle{Научная новизна}
    \small
    \begin{columns}[T,onlytextwidth]
        \begin{column}{0.48\textwidth}
            \begin{block}{Модельное многообразие}
                Предложена методика построения и исследования регулярных ветвей решений задачи встречи методом дискретного продолжения.
            \end{block}
            \begin{block}{Конечные формулы}
                Получены конечные формулы для \(T\), \(\Psi\), \(J\) и \(\mathbf z_0\) как функций параметров модельной задачи.
            \end{block}
        \end{column}
        \begin{column}{0.48\textwidth}
            \begin{block}{Линейные \(R\)-преобразования}
                Предложены преобразования начальных значений сопряжённых переменных для переноса модельной экстремали к другой краевой задаче.
            \end{block}
            \begin{block}{Алгоритм}
                Разработан и программно реализован алгоритм выбора начального приближения для прикладных классов перелётов.
            \end{block}
        \end{column}
    \end{columns}
\end{frame}

\section{Модельное многообразие}

\begin{frame}[t]
    \frametitle{ИРТ как модельная краевая задача и параметры}
    \begin{columns}[T,onlytextwidth]
        \begin{column}{0.48\textwidth}
            \tiny
            \[
            \left\{
            \begin{aligned}
            \dot{\mathbf r} &= \mathbf v,\\
            \dot{\mathbf v} &= -\frac{\mathbf r}{r^3}+\mathbf a,
            \end{aligned}
            \right.
            \]
            \[
            J=\frac{1}{2}\int\limits_{t_0}^{t_f}\|\mathbf a(t)\|^2\,dt \to \min.
            \]
            \[
            \mathbf z_0=\left(\mathbf p_{r0},\mathbf p_{v0}\right),
            \qquad
            \mathbf a=\mathbf p_v.
            \]
            \[
            (\lambda,\kappa)
            \longrightarrow
            T,\ \Psi,\ J,\ \mathbf z_0 .
            \]
            \begin{itemize}
                \item \(\kappa\) задаёт отношение радиусов модельных круговых орбит.
                \item \(\lambda\) задаёт витковую дальность.
                \item \(\Phi\) связывает выбранную ветвь с датой старта.
                \item ИРТ используется как регулярная модельная КЗ для построения начального приближения.
            \end{itemize}
        \end{column}
        \begin{column}{0.48\textwidth}
            \centering
            \includegraphics[width=\linewidth,height=0.66\textheight,keepaspectratio]{virtual_planet}
        \end{column}
    \end{columns}
\end{frame}

\begin{frame}[t]
    \frametitle{Модельное многообразие и дискретное продолжение}
    \begin{columns}[T,onlytextwidth]
        \begin{column}{0.58\textwidth}
            \centering
            \includegraphics[height=0.48\textheight]{continuation_paths}
        \end{column}
        \begin{column}{0.38\textwidth}
            \tiny
            \[
            \mathbf z(\boldsymbol\xi)\in\mathbb R^6,
            \qquad
            \mathbf c(\boldsymbol\xi)\in\mathbb R^n.
            \]
            \[
            M:D_{\mathrm{reg}}\to\mathbb R^{6+n},
            \qquad
            M(\boldsymbol\xi)=
            \bigl(\mathbf z(\boldsymbol\xi),\mathbf c(\boldsymbol\xi)\bigr).
            \]
            \begin{itemize}
                \item Модельное многообразие понимается как образ отображения, построенного на одной регулярной ветви решений.
                \item Базовое решение используется как исходная точка для продолжения по сетке параметров.
                \item Соседние решения дают начальный вектор для следующей КЗ.
                \item Вычислительные барьеры отделяют границы устойчивого локального продолжения.
            \end{itemize}
        \end{column}
    \end{columns}
    \begin{block}{Новый результат}
        \scriptsize Разработана методика построения регулярного модельного многообразия методом дискретного продолжения.
    \end{block}
\end{frame}

\begin{frame}
    \frametitle{От многообразия к конечным формулам}
    \small
    \[
    M(\boldsymbol\xi)
    =
    \bigl(\mathbf z(\boldsymbol\xi),\mathbf c(\boldsymbol\xi)\bigr),
    \qquad
    \mathbf c=(\Phi,\lambda,T,J).
    \]
    \[
    M(\boldsymbol\xi)
    \longrightarrow
    \widehat T,\ \widehat\Psi,\ \widehat J,\ \widehat{\mathbf z}_0.
    \]
    \begin{itemize}
        \item Численные решения анализируются в выбранных параметрических представлениях.
        \item Для регулярных зависимостей строятся непосредственные аппроксимации.
        \item Для более сложных зависимостей используется иерархический подход с коэффициент-функциями.
        \item Символьная регрессия используется как инструмент получения конечных формул.
        \item Качество формул проверяется метриками ошибок и краевыми невязками после интегрирования расширенной системы.
    \end{itemize}
    \begin{block}{Новый результат}
        \scriptsize Задана методика получения интерпретируемых конечных формул для характеристик экстремалей и начальных значений сопряжённых переменных.
    \end{block}
\end{frame}

\begin{frame}[t]
    \frametitle{КЛО и точки оптимальной дальности}
    \begin{columns}[T,onlytextwidth]
        \begin{column}{0.58\textwidth}
            \centering
            \includegraphics[height=0.48\textheight]{Isolines}
        \end{column}
        \begin{column}{0.38\textwidth}
            \scriptsize
            \[
            \Phi=\mathrm{const}
            \quad\cap\quad
            \mathrm{КЛО}
            \quad\Rightarrow\quad
            \lambda^\star.
            \]
            \[
            \varphi=\Phi+\frac{T}{\kappa^{3/2}}.
            \]
            \begin{itemize}
                \item Дата старта задаёт начальную синодическую фазу.
                \item КЛО строится как отдельная ветвь решений.
                \item Пересечения КЛО с изолиниями фазы образуют точки оптимальной дальности.
                \item Точка оптимальной дальности задаёт \(\lambda^\star\).
            \end{itemize}
        \end{column}
    \end{columns}
    \begin{block}{Новый результат}
        \scriptsize Выделены КЛО и точки оптимальной дальности как механизм выбора ветви по фиксированной начальной синодической фазе.
    \end{block}
\end{frame}
